Can daily data distinguish jumps from asymmetric heavy tails, and how should the empirical ranking be read? To answer these questions we simulate from five data-generating processes (DGPs) whose parameters equal the full-sample estimates in Table~\ref{tab:insample}: Gaussian, i.i.d.\ Merton, GARCH-$t$, GARCH-skew-$t$ and GARCH-Merton. In each of \mcReps{} replications per DGP we simulate \mcTrain{}~+~\mcTest{} returns, estimate all eight models on the first \mcTrain{} observations (the length of the rolling window) and evaluate one-step-ahead forecasts on the last \mcTest{}, with GARCH variances filtered from past simulated returns. Table~\ref{tab:mc} reports mean negative log scores, 99\% VaR exceedance rates and the share of replications in which each model has the best score.

\begin{table}[t]
\centering\small
\caption{Monte Carlo results: \mcReps{} replications per DGP, \mcTrain{} estimation and \mcTest{} evaluation observations. First row of each block: mean negative log score (bold: lowest). Second row: mean 99\% VaR exceedance rate in \% (nominal 1\%). Third row: share of replications in which the model has the lowest log score.}\label{tab:mc}
\resizebox{\linewidth}{!}{\begin{tabular}{lrrrrrrrr}
\toprule
DGP & Gauss & $t$ & Merton & Kou & G-$t$ & G-skew-$t$ & G-Merton & G-Kou \\
\midrule
Gaussian & \textbf{1.0477} & 1.0478 & 1.0481 & 1.0479 & 1.0480 & 1.0484 & 1.0482 & 1.0481 \\
\quad 99\% exc. (\%) & 1.00 & 0.97 & 1.00 & 1.00 & 0.97 & 0.97 & 1.01 & 1.01 \\
\quad Best NLS (\%) & 24 & 12 & 16 & 2 & 14 & 16 & 8 & 10 \\
\midrule
Merton & 1.0232 & 0.9474 & \textbf{0.9420} & 0.9424 & 0.9477 & 0.9452 & 0.9423 & 0.9427 \\
\quad 99\% exc. (\%) & 2.50 & 1.55 & 1.01 & 1.00 & 1.54 & 1.13 & 1.01 & 1.00 \\
\quad Best NLS (\%) & 0 & 4 & 30 & 16 & 3 & 7 & 26 & 13 \\
\midrule
GARCH-t & 1.1707 & 0.9765 & 0.9902 & 0.9829 & \textbf{0.8644} & 0.8648 & 0.8721 & 0.8715 \\
\quad 99\% exc. (\%) & 1.96 & 1.27 & 1.19 & 1.19 & 0.98 & 1.00 & 1.14 & 1.13 \\
\quad Best NLS (\%) & 0 & 0 & 0 & 0 & 50 & 34 & 8 & 8 \\
\midrule
GARCH-skew-t & 1.1582 & 0.9586 & 0.9732 & 0.9632 & 0.8472 & \textbf{0.8412} & 0.8519 & 0.8497 \\
\quad 99\% exc. (\%) & 2.38 & 1.69 & 1.35 & 1.33 & 1.51 & 1.04 & 1.23 & 1.18 \\
\quad Best NLS (\%) & 0 & 0 & 0 & 0 & 11 & 77 & 6 & 6 \\
\midrule
GARCH-Merton & 0.9890 & 0.9415 & 0.9403 & 0.9403 & 0.8704 & 0.8682 & 0.8637 & \textbf{0.8635} \\
\quad 99\% exc. (\%) & 2.34 & 1.58 & 1.24 & 1.15 & 1.51 & 1.26 & 1.08 & 1.12 \\
\quad Best NLS (\%) & 0 & 0 & 0 & 0 & 4 & 8 & 44 & 44 \\
\bottomrule
\end{tabular}
}
\end{table}

We draw three conclusions. First, the mechanism is identifiable at this sample size. The Bayesian information criterion selects the true model in \mcBicGarchMertonGarchMerton\% of GARCH-Merton samples, \mcBicGarchSkewTGarchSkewT\% of GARCH-skew-$t$ samples and \mcBicGarchTGarchT\% of GARCH-$t$ samples. Out of sample, the correctly specified class also has the lowest mean score for every DGP; under the GARCH-Merton DGP the two jump-GARCH models are essentially tied, because a Kou mixture approximates a Merton mixture well.

Second, the simulations show how to read the empirical evidence. When data come from the GARCH-Merton process, the jump models beat GARCH-$t$ (mean scores \mcNlsGarchMertonGarchMerton{} and \mcNlsGarchMertonGarchT) and the skewed-$t$ model lies in between (\mcNlsGarchMertonGarchSkewT). When data come from the GARCH-skew-$t$ process, the order is skewed-$t$ first (\mcNlsGarchSkewTGarchSkewT), GARCH-$t$ second (\mcNlsGarchSkewTGarchT) and the jump models last (\mcNlsGarchSkewTGarchMerton{} and \mcNlsGarchSkewTGarchKou). The empirical ranking in Table~\ref{tab:compare}, with the skewed-$t$ model better and the jump models worse than GARCH-$t$, matches the second pattern and not the first. Coverage agrees as well: under the skewed-$t$ DGP the 99\% VaR of the GARCH-$t$ model is exceeded on \mcExcGarchSkewTGarchT\% of days, close to the \rateNNGarchT\% observed for the peso, and the jump models reduce this to \mcExcGarchSkewTGarchKou--\mcExcGarchSkewTGarchMerton\%. The partial repair of coverage by jump models is therefore exactly what one would expect if the peso had asymmetric tails and no jumps at all.

Third, even when jumps are present, their parameters are estimated imprecisely. Under the GARCH-Merton DGP, with true intensity $\lambda=\mcTrueLambda$, the root mean squared error of $\hat\lambda$ is \mcRmseLambda{} and 90\% of the estimates lie between \mcLoLambda{} and \mcHiLambda; for the mean jump size ($\mu_J=\mcTrueMuj$) the root mean squared error is \mcRmseMuj. The instability of the rolling jump estimates in Figure~\ref{fig:params} is of the size that sampling error alone would produce. Finally, when data are i.i.d.\ (Gaussian or Merton DGPs), the GARCH models lose very little relative to the true model, so allowing for volatility clustering carries almost no cost.
